\BeleskeID{09_2-s052}
% element: 09_2-s052:e04
Lečap može biti destruktivan, ali ograničen strujni kapacitet izvora ponekad sprečava da disipacija dostigne destruktivnu vrednost. Tada pad napona napajanja ipak onemogućava ispravno funkcionisanje. Za prekid već uspostavljenog lečapa uklanja se napajanje.

Pojava je česta na delovima povezanim sa spoljnim svetom: pri uspostavljanju ili nestajanju napona izlazni priključci mogu imati nagle promene. Povećanjem vremena uspona i pada smanjuje se brzina promene i pobuda kroz parazitne kapacitivnosti. Druga tehnika okružuje oblasti zaštitnim prstenovima povezanim na odgovarajući potencijal napajanja ili mase. Time se smanjuje međusobni uticaj parazitnih tranzistora.
